\documentclass[9pt,twocolumn]{article}
\usepackage[letterpaper,margin=0.55in,top=0.48in,bottom=0.52in]{geometry}
\usepackage[T1]{fontenc}\usepackage{newtxtext,newtxmath}\usepackage{microtype}
\usepackage{amsmath,mathtools}\usepackage{booktabs,tabularx,array}\usepackage[normalem]{ulem}\usepackage{cancel}\usepackage{enumitem,xcolor,ragged2e}
\usepackage{graphicx}
\setlength{\columnsep}{0.24in}\setlength{\parindent}{0pt}\setlength{\parskip}{2.2pt}
\setlist{nosep,leftmargin=*}\renewcommand{\arraystretch}{1.10}
\newcommand{\oldclaim}[1]{\sout{#1}}
\newcommand{\claimbox}[1]{\begin{center}\fbox{\parbox{0.94\columnwidth}{\centering\scriptsize #1}}\end{center}}
\newcommand{\eclass}[1]{\textsf{[#1]}}
\DeclareFontFamilySubstitution{TS1}{ntxtlf}{cmr}
% \let binds the CURRENT meaning of \_ ; \newcommand would create a macro that
% re-expands to \_ and loops until TeX runs out of input stack.
\let\CIGunderscore\_
\renewcommand{\_}{\CIGunderscore\allowbreak}
\begin{document}
\sloppy
\twocolumn[{\begin{center}
{\LARGE\bfseries Report 2 - Functional Response Geometry}\\[2pt]
{\large\itshape Append-only update E: the extremal-geometry effect, and the covariance design module}\\[2pt]
{\small CIG-AMF Scientific Evolution Report -- 11 September (backlog entry)}
\end{center}\vspace{2mm}}]

\textbf{Append-only rule.} Sections 1--70 are an immutable prefix. This addendum begins at Section 71.

\section*{71. Extremal geometry, outer gap and noise held fixed}
Holding the outer Top-$K$ gap and the noise level constant and changing only the \emph{geometry of the extrema}, the anytime variant's sample cost falls from $9\,167$ to $5\,700$, and it is better on \textbf{all twelve} paired configurations. Because the gap and noise are fixed by construction, the change cannot be attributed to instance difficulty in the usual sense; it isolates the geometry of where the extrema sit within each relation.
\claimbox{\eclass{DEVELOPMENT EXPERIMENT}. Twelve paired configurations is a design, not a confirmatory study, and the generator still fixes extremum positions in the allocation branch --- so this measures a contrast between two designed geometries, not a law. It does, however, identify the quantity a maintenance-cost theorem would have to contain: not total drift, but drift that moves comparisons across the Top-$K$ boundary.}

\section*{72. Module recorded}
\texttt{Lean\allowbreak Functional\allowbreak Covariance\allowbreak DesignV1} joins the chain's Lean material alongside \texttt{LeanFunctionalResidualMaintenanceV2}, whose \texttt{selective\_refresh\_safety} takes both the coverage hypothesis \texttt{hCoverage} and the future-error bound \texttt{hFuture} as \emph{premises}. That remains the chain's central theoretical gap: how coverage and the drift bound are \emph{learned}, and at what cost, is unproved.

\section*{73. Standing collisions, restated so they are in the ledger}
Verified directly, not from memory: structured best-arm identification with micro-observables mapping to per-action scores (Huang, Ajallooeian, Szepesv\'ari, M\"uller, ALT 2017, with LUCB-micro and matching bounds); composable safe zones with a proved composition theorem (Garofalakis and Samoladas, ICDT 2017); best-arm identification with possibly biased offline data (Yang, Tan, Cheung, UAI 2025). The chain's two-tier structure, its safe-zone-style selective refresh, and its reuse of stale data each sit against one of these. What remains open to this chain is the \emph{maintenance cost law} under drift with changing extrema, measured against restart and safe-zone monitoring at equal cost --- not the mechanisms themselves.
\end{document}
